% ЕДИНЫЙ ШАБЛОН. Компилятор: XeLaTeX.
% Сначала измените реквизиты ниже, затем замените примеры своим содержанием.
\documentclass{mietlab}

\LabSetup{
  number=1,
  title={«Тест-дизайн. Техники тест-дизайна»},
  student={Варлапов Артем},
  group={П-42},
  discipline={Основы программной инженерии},
  city={Москва},
  year={2026}
}

\begin{document}
\MakeLabTitle

% Заголовок и текст в одном абзаце, как «Цель работы» в первой лабе.
\LabLabel{Цель работы}{Изучить основные техники тест-дизайна.}

\LabHeading{Задание}
Здесь укажите вариант, условие задания и исходные данные.
Абзацы выравниваются по ширине, без красной строки; межстрочный интервал — 1,5.

\LabHeading{Выполнение работы}
\LabHeading{1. Требования к поисковой форме}
Вариант: \url{http://www.rambler.ru/}. Для открытия формы поиска по интернету
используется ссылка «Поиск» на главной странице Рамблера. Она ведёт на страницу
\url{https://r0.ru/}.

% Таблица без названия, как в первой лабе. Заголовки полужирные,
% линии 0.5 bp, текст сверху, боковые поля ячеек 5 bp.
% Доли первых двух столбцов взяты из исходной таблицы: 105/482 и 100/482.
% У последнего столбца вычитается толщина крайней линии сетки.
\begin{labtable}{|L{0.21784\linewidth}|L{0.20747\linewidth}|L{\dimexpr0.57469\linewidth-\arrayrulewidth\relax}|}
\hline
\LabTableHead{Элемент} & \LabTableHead{Тип} & \LabTableHead{Требования} \\\hline
\endfirsthead
\hline
\LabTableHead{Элемент} & \LabTableHead{Тип} & \LabTableHead{Требования} \\\hline
\endhead
Поле поиска & Текстовое поле &
1. Принимает буквы, цифры, пробелы и специальные символы.\par
2. Допускает запрос из одного непробельного символа.\par
3. Позволяет изменять и удалять введённый текст. \\\hline
Проверка ввода & Обработка данных &
Пустой запрос и строка из пробелов не запускают поиск.
После удаления текста прежний запрос не отправляется. \\\hline
Кнопка поиска & Кнопка &
Отправляет текущее непустое значение поля и открывает страницу результатов. \\\hline
Клавиша Enter & Действие клавиатуры &
При фокусе в поле запускает поиск по текущему запросу. \\\hline
Страница результатов & Веб-страница &
Отображает введённый запрос и найденные результаты либо сообщение об их отсутствии.
Кириллица отображается без искажений. \\\hline
\end{labtable}

% Удалите принудительный разрыв, если раздел должен продолжаться на странице.
\LabNewPage
\LabHeading{2. Таблица тестовых данных}
Данные подобраны с помощью эквивалентного разбиения (EP), анализа граничных
значений (BVA) и предугадывания ошибок (EG). OK обозначает допустимые данные,
NOK — недопустимые.

\begin{labtable}{|L{0.13071\linewidth}|L{0.14938\linewidth}|L{0.30083\linewidth}|L{\dimexpr0.41908\linewidth-\arrayrulewidth\relax}|}
\hline
\LabTableHead{Поле} & \LabTableHead{OK/NOK} & \LabTableHead{Значение} & \LabTableHead{Комментарий} \\\hline
\endfirsthead
\hline
\LabTableHead{Поле} & \LabTableHead{OK/NOK} & \LabTableHead{Значение} & \LabTableHead{Комментарий} \\\hline
\endhead
Поиск & OK & Rambler & Латиница (EP). \\\hline
Поиск & OK & тестирование ПО & Кириллица и пробел между словами (EP). \\\hline
Поиск & OK & 2026 & Запрос из цифр (EP). \\\hline
Поиск & OK & C++ & Буква и специальные символы (EP, EG). \\\hline
Поиск & NOK & Пустая строка & Ноль символов, значение ниже нижней границы (EP, BVA). \\\hline
\end{labtable}

\LabHeading{3. Тест-кейс}
\LabLabel{Предусловия}{Открыта форма поиска; поле пустое.}
\LabLabel{Данные}{Rambler.}

\begin{labtable}{|L{0.45643\linewidth}|L{\dimexpr0.54357\linewidth-\arrayrulewidth\relax}|}
\hline
\LabTableHead{Действие} & \LabTableHead{Ожидаемый результат} \\\hline
\endfirsthead
\hline
\LabTableHead{Действие} & \LabTableHead{Ожидаемый результат} \\\hline
\endhead
1. Ввести выбранное значение. & В поле отображается введённый запрос. \\\hline
2. Нажать кнопку поиска. & Открываются результаты для введённого запроса
либо сообщение об их отсутствии. Запрос отображается без искажений. \\\hline
\end{labtable}

\LabNewPage
\LabHeading{4. Пример программного кода}
% Во втором PDF код был 8.5 pt. Здесь он намеренно 13 pt,
% Times New Roman, интервал 1.5, согласно вашему требованию.
\begin{labcode}[language=Java]
public class Book {
    private String title;

    public Book(String title) {
        this.title = title;
    }

    public String getTitle() {
        return title;
    }
}
\end{labcode}

\LabHeading{5. Рисунок}
Результат выполнения тестов представлен на рисунке~\ref{fig:result}.
% Замените картинку и текст подписи. Четвёртый аргумент — уникальная метка.
% Это скриншот из второго пользовательского PDF; его содержимое не меняется.
\LabFigure{assets/example-tests.png}{Результат выполнения тестов}{fig:result}

\LabHeading{Вывод}
Здесь сформулируйте результат работы: что реализовано, что проверено,
какие выводы сделаны. Этот текст — заполняемый пример для новой лабораторной работы.

\end{document}
